% !TEX root = ../../main.tex

\section{Background and Motivation}
\label{sec:carp:bg}

Data generated by many scientific applications consists of simulated entities such as particles and
mesh cells, each storing several attributes (e.g. energy, velocity) into a \textit{record}.
This data is written out at fixed intervals and not mutated once written, making bulk indexing
feasible.
The data is then analyzed using point queries (e.g., retrieving particles by ID) or range queries
(e.g., retrieving particles with energy in a certain range).
Range queries are needed for interacting with continuous attributes and for filtering data using
ranges of relevant values (e.g. energy bands, temperature and blast wave fronts).

An efficient range-query index must locate ranges of values quickly.
Sorted indexes do this by creating indexed attribute (henceforth referred to as \emph{key})
locality in the data layout, so that the keys for a range can be located in one or few data
partitions.
Clustered indexes further improve query performance by storing all records with their corresponding
keys.
This allows for data for a query to be retrieved efficiently via fewer, contiguous reads.
In contrast, auxiliary indexes need multiple random reads to retrieve all attributes of interest in
the range.
In general, a database table can have only one clustered index.
The clustered index is typically reserved for the attribute most commonly used in
queries~\cite{MikeR2023:ClusteredNonclusteredIndexes,Ray2023:ClusteredIndexes}.

\newcolumntype{C}{>{\centering\arraybackslash}m{3.0cm}}
\newcolumntype{L}{>{\raggedleft\arraybackslash}m{4.85cm}}
\renewcommand{\arraystretch}{1.2}

\begin{table}[t]
  \caption{Range query indexing approaches.
  }
  \label{tab:carp:comp}
  \centering
  \begin{tabular}{LCCC}
    \toprule
    Approach                      & In-situ    & Efficient Indexing & Efficient Querying \\
    \midrule
    Bulk Sorting\- (Clustered)    & $\times$   & $\times$           & \checkmark         \\
    Bulk Sorting (Auxiliary)      & $\times$   & $\times$           & $\times$           \\
    Bitmap Indexes (FastQuery)    & $\times$   & $\times$           & $\times$           \\
    DB Indexes (LSM-Tree, B-Tree) & \checkmark & $\times$           & \checkmark         \\
    DeltaFS                       & \checkmark & $\checkmark$       & $\times$           \\
    \textbf{CARP}                 & \checkmark & \checkmark         & \checkmark         \\
    \bottomrule
  \end{tabular}
\end{table}


\cref{tab:carp:comp} shows that no current approach enables both in-situ efficient indexing and
efficient querying for ranges.
% \vspace{-2\baselineskip}
We summarize the characteristics of the indexing approaches below.
We discuss their \emph{index construction} performance and break down their query performance into
\emph{index lookup} and \emph{data retrieval}.

\textbf{Index Construction.
}
Bulk parallel/distributed sorting~\cite{Sunda2013:HykSort, Dong2016, Rasmu2011:TritonSort}
approaches support range queries but require post-processing.
These can be used to build either clustered or auxiliary sorted indexes.
FastQuery~\cite{Chou2011:FastQuery}, on the other hand, builds auxiliary bitmap structures where
keys are sorted and stored separately with a pointer to their records.
Clustered indexes enable more efficient queries but require more I/O to reorder and write the
entire dataset.
FastQuery does not reorder data, but results in data scans, computing large bitmap vectors, and
writing auxiliary structures to disk.

To avoid expensive post-processing, techniques have been proposed to intercept application writes
and index them in-situ before data is written to disk.
DeltaFS \cite{Zheng2018:DeltaFS} shuffles data into hash-based partitions and requires no
post-processing, but it can only support point queries on the partitioned data.
Distributed databases maintain online indexes but offer lower write efficiency as they periodically
reorganize on-disk data \cite{Raju2017:PebblesDB}.
Databases are designed for mixed workloads with updates, while lightweight indexes are better
suited to scientific data.

\looseness=-1
\textbf{Index Lookup and Data Retrieval.
}
Efficient range queries require quickly finding matching records using the index and retrieving the
records from storage using large sequential reads.
Only sorted and clustered indexes provide both of these properties.
Auxiliary indexes can return the offsets at which matching rows are present, but retrieving those
records requires multiple random reads.
Additionally, FastQuery queries require provisioning enough compute nodes to fit large index
structures in memory, and queries can only be processed once these indexes are loaded.

\looseness=-1
With CARP, we aim to pair the query latency of a sorted and clustered index with the reduced
overhead of an in-situ approach.
% We discuss the challenges of meeting these goals in the next section.
We discuss the resulting challenges in the next section.
